
\documentclass{../notatki}

\title{Matematyczne fundamenty informatyki}

\begin{document}

\section{Analiza Matematyczna}

\subsection{Pochodna funkcji złożonej}
$$
h = f \circ g \quad h'(x) = f'(g(x)) \cdot g'(x)
$$

\begin{examplebox}
  Pochodną funkcji $f(x) = (x^9 + x)^9$ jest:
  $$
  f'(x) = 9(x^9 + x)^8 \cdot (9x^8 + 1)
  $$
\end{examplebox}

\subsection{Jakobian}
$$
f: \mathbb{R}^n \to \mathbb{R}^m \quad f(x) = (f_1(x), \dots, f_m(x))
$$
Dla powyższej funkcji, macierz Jakobiego to:
$$
J_{ij} = \frac{\partial f_i}{\partial x_j} \quad J_f =
\begin{bmatrix}
  \frac{\partial f_1}{\partial x_1} & \dots & \frac{\partial
  f_1}{\partial x_n} \\
  \vdots & \ddots & \vdots \\
  \frac{\partial f_m}{\partial x_1} & \dots & \frac{\partial f_m}{\partial x_n}
\end{bmatrix}
$$
Jakobian, to wyznacznik macierzy Jakobiego.

\begin{examplebox}
  $$
  f(x, y, z) = (x, y, z) \quad f_1(x, y, z) = x \quad f_2(x, y, z) =
  y \quad f_3(x, y, z) = z
  $$
  $$
  J_f =
  \begin{bmatrix}
    1 & 0 & 0 \\
    0 & 1 & 0 \\
    0 & 0 & 1
  \end{bmatrix}
  $$
  $$
  \det(J_f) = 1
  $$
\end{examplebox}

\section{Rachunek Prawdopodobieństwa}

Gęstością pewnego zdarzenia nazywamy funkcję $f: \Omega \to [0, 1]$, która
dla każdego $x \in \Omega$ zwraca prawdopodobieństwo zdarzenia $x$.
Wymogiem formalnym gęstości, jest to, że się sumuje do 1:
$$
\sum_{x \in \Omega} f(x) = 1 \quad \int_{\Omega} f(x) dx = 1
$$

\begin{examplebox}
  Na przykład, dla zmiennej losowej $X$, określonej przez wynik losowego
  rzutu koścą, prawdopodobieństwo zdarzenia $X = k$ wynosi
  $\frac{1}{6}$, gdzie $k$ to wynik rzutu. Stąd też, $f(x) =
  \frac{1}{6}$ dla $x = k$.
\end{examplebox}

Dystrybuanta, to funkcja $F: \Omega \to [0, 1]$, ściśle związana z gęstością:
$$
F(x) = P(X \leq x) = \int_{-\infty}^x f(t) dt
$$
W pewnym sensie, reprezentuje prawdopodobieństwo, że zmienna losowa
$X$ ma wartość mniejszą lub równą $x$.

\begin{examplebox}
  Wracając do naszego przykładu z kością, dystrybuanta $F(x)$ dla $x = k$ wynosi
  $$
  F(x) = P(X \leq x) = \int_{-\infty}^x f(t) dt = \int_{-\infty}^k
  \frac{1}{6} dt = \frac{k}{6}
  $$
  W tym przypadku, $F(x)$ jest funkcją skokową, która przyjmuje
  wartość 0 dla $x < 1$, 1/6 dla $x = 1$, 2/6 dla $x = 2$, itd.
\end{examplebox}
\textbf{Dystrybuanta nie musi się sumować do 1}, wręcz przeciwnie, ale musi być
niemalejąca, prawostronnie ciągła (lub lewostronnie, ale nie może mieć dziur).
$$
\lim_{x \to \infty} F(x) = 1 \quad \lim_{x \to -\infty} F(x) = 0
$$

\begin{examplebox}
  Poniższa funkcja $F(x)$ jest dystrybuantą:
  $$
  F(x) =
  \begin{cases}
    0 & x < 1 \\
    1 - \frac{1}{x^2} & x \geq 1
  \end{cases}
  $$
  Ponieważ jest prawostronnie ciągla, niemalejąca i ograniczona przez 1.

  \centering
  \begin{tikzpicture}
    \begin{axis}[
        axis lines=middle,
        xlabel={$x$},
        ylabel={$F(x)$},
        ymin=-0.1, ymax=1.1,
        xmin=0, xmax=4,
        samples=200,
        domain=0:4,
        width=10cm,
        height=6cm,
        xtick={1},
        ytick={0,1},
        legend pos=north east,
        clip=false
      ]
      % F(x) = 0 for x < 1
      \addplot[blue, thick, domain=0:1, samples=2] {0};
      % F(x) = 1 - 1/x^2 for x >= 1
      \addplot[blue, thick, domain=1:4] {1 - 1/(x^2)};
    \end{axis}
  \end{tikzpicture}
\end{examplebox}

\section{Algebra Liniowa}

Wektor to element przestrzeni wektorowej $V$, czyli $x \in V$. Wektory są
ortogonalne, jeśli ich iloczyn skalarny jest równy 0, czyli $x \cdot y = 0$ dla
dowolnych $x, y \in V$.

\begin{examplebox}
  Wektory $(1, 1, 1)$ i $(-2, 1, 1)$ są ortogonalne, ponieważ ich
  iloczyn wynosi:
  $$
  \begin{bmatrix}
    1 & 1 & 1
  \end{bmatrix}
  \begin{bmatrix}
    -2 \\ 1 \\ 1
  \end{bmatrix}
  = 0
  $$
\end{examplebox}

Normą indukowaną w $\mathbb{R}^n$ jest funkcja $||x|| = \sqrt{x^T
\cdot x} = \sqrt{\braket{x}{x}}$.

\subsection{Wyznacznik}

$$
\det A = a_{11} \det A_{11} - a_{12} \det A_{12} + \dots + (-1)^{n-1}
a_{nn} \det A_{nn}
$$
gdzie $A_{ij}$ to macierz $A$ z usuniętym wierszem $i$ i kolumną $j$.

\begin{examplebox}
  Dla macierzy $A$ rozmiaru $3 \times 3$:
  $$
  \det A = a_{11}
  \begin{vmatrix}
    a_{22} & a_{23} \\
    a_{32} & a_{33}
  \end{vmatrix} - a_{12}
  \begin{vmatrix}
    a_{21} & a_{23} \\
    a_{31} & a_{33}
  \end{vmatrix} + a_{13}
  \begin{vmatrix}
    a_{21} & a_{22} \\
    a_{31} & a_{32}
  \end{vmatrix}
  $$

  Dla macierzy $A$ rozmiaru $4 \times 4$:
  $$
  \det A = a_{11}
  \begin{vmatrix}
    a_{22} & a_{23} & a_{24} \\
    a_{32} & a_{33} & a_{34} \\
    a_{42} & a_{43} & a_{44}
  \end{vmatrix} - a_{12}
  \begin{vmatrix}
    a_{21} & a_{23} & a_{24} \\
    a_{31} & a_{33} & a_{34} \\
    a_{41} & a_{43} & a_{44}
  \end{vmatrix} + a_{13}
  \begin{vmatrix}
    a_{21} & a_{22} & a_{24} \\
    a_{31} & a_{32} & a_{34} \\
    a_{41} & a_{42} & a_{44}
  \end{vmatrix} - a_{14}
  \begin{vmatrix}
    a_{21} & a_{22} & a_{23} \\
    a_{31} & a_{32} & a_{33} \\
    a_{41} & a_{42} & a_{43}
  \end{vmatrix}
  $$

\end{examplebox}

\subsection{Odwzorowania liniowe}

Odwzorowanie liniowe $T: V \to U$, to
macierzowa transformacja
$$
T(x) = Ax
$$
gdzie $A$ to macierz $\dim V \times \dim U$, a $x$ to wektor $\dim
V$-elementowy.

$$
\ker T = \{x : T(x) = 0\} \quad \Im T = \{T(x) : x \in V\}
$$

Zbiór wektorów jest liniowo niezależny, gdy nie istnieje żadna
kombinacja liniowa tych wektorów, która daje wynik równy 0.
$$
\alpha_1 v_1 + \alpha_2 v_2 + \dots + \alpha_n v_n = 0
$$

Ranga macierzy $A$ to maksymalna liczba liniowo niezależnych wierszy
(lub kolumn) macierzy $A$, czyli $rank(A) = \dim \Im T$. Formalnie;
wymiar przestrzeni rozpiętej przez jej wiersze. Jeśli ranga jest
równa zero, to macierz wysyła wektory do zera, czyli $\ker T = V$.

Operator liniowy jest różnowartościowy wtedy i tylko wtedy, gdy $\ker
T = \{0\}$. Macierz przekształcenia musi wtedy też mieć rangę $n$, gdzie $n$ to
liczba wierszy macierzy $A$.

$$
\dim \mathbb{R}^n = n \quad \dim V = \dim \Im T + \dim \ker T
$$

\begin{examplebox}
  Jeśli wiemy, że $T: \mathbb{R}^4 \to \mathbb{R}^3$, to wiemy, że macierz
  $T$ ma rozmiar $3 \times 4$.
  $$
  T(x) =
  \begin{bmatrix}
    a_{11} & a_{12} & a_{13} & a_{14} \\
    a_{21} & a_{22} & a_{23} & a_{24} \\
    a_{31} & a_{32} & a_{33} & a_{34}
  \end{bmatrix}
  \begin{bmatrix}
    x_1 \\
    x_2 \\
    x_3 \\
    x_4
  \end{bmatrix}
  =
  \begin{bmatrix}
    a_{11}x_1 + a_{12}x_2 + a_{13}x_3 + a_{14}x_4 \\
    a_{21}x_1 + a_{22}x_2 + a_{23}x_3 + a_{24}x_4 \\
    a_{31}x_1 + a_{32}x_2 + a_{33}x_3 + a_{34}x_4
  \end{bmatrix}
  $$

  $T$, może mieć nietrywialne $\ker T \neq \{0\}$, może mieć rangę w
  $[0, 4]$ bo ma 4 kolumny, lecz nie może być różnowartościowe, ponieważ:
  $$
  \dim V = \dim \Im T + \dim \ker T \Rightarrow 4 = 3 + \dim \ker T
  \Rightarrow \dim \ker T = 1
  $$
  Więc $\ker T \neq \{0\}$, więc $T$ nie jest różnowartościowy.
\end{examplebox}

Aby macierz była odwracalna, to musi być kwadratowa, czyli $n \times
n$, gdzie $n$ to liczba wierszy, jej ranga musi być maksylana $\dim \Im T = n$,
więc $\dim \ker T = 0$, czyli $\ker T = \{0\}$, i jej wyznacznik musi
być różny od zera, czyli $\det T \neq 0$.

\subsection{Wartości własne}

Dla równania:
$$
A \ket{\psi} = \lambda \ket{\psi}
$$
$\ket{\psi}$ jest wektorem własnym macierzy $A$, a $\lambda$ jest jej
wartością własną - odpowiadającą wektorowi własnemu $\ket{\psi}$.
$$
\det(A - \lambda I) = 0 \quad (A - \lambda I) \ket{\psi} = 0
$$

Wielomian charakterystyczny macierzy to:
$$
p_A(t) = (t - \lambda_1)^{k_1} (t - \lambda_2)^{k_2} \cdots (t -
\lambda_n)^{k_n} = 0 \quad p_A(A) = 0
$$
jest to zazwyczaj pochodna poszukiwania wartości i wektorów własnych.

\subsection{Rozkłady}

Każdą rzeczywistą, symetryczną macierz dodatnio określonej formy kwadratowej
można zapisać jako iloczyn macierzy dolnotrójkątnej i jej transpozycji.
Nazywamy to rozkładem Choleskiego:
$$
A = LL^T
$$

\section{Statystyka}

Estymator, to funkcja $\hat{\theta}(X)$, która na podstawie próbki $X
\subset U$, gdzie $U$ to cała populacja, estymuje wartość parametru $\theta(U)$.

Estymator jest nieobiążony, jeśli:
$$
\mathbb{E}[\hat{\theta}(X)] = \theta(U)
$$
czyli wartość oczekiwana estymatora jest równa wartości parametru. Mówimy wtedy,
że $\hat{\theta}(X)$ jest estymatorem największej wiarygodności (ENW).

\begin{examplebox}
  Dla próbki $X \subset U$, gdzie $U$ to cała populacja, estymator
  średniej $\hat{\mu}(X) = \overline{X}$, jest nieobciążonym
  estymatorem średniej $\mu(U)$, ponieważ:
  $$
  \mathbb{E}[\hat{\mu}(X)] = \mu(U)
  $$
  Intuicyjnie, estymator jest nieobciążony, ponieważ dla odpowiednio dużej
  próbki $X$, średnia z próbki będzie bliska wartości parametru $\mu(U)$.
\end{examplebox}

\subsection{Charakterystyka próbek losowych}

Dla zmiennej losowej $X$,
\begin{itemize}
  \item $\mathbb{E}[X] = \mu$ - wartość oczekiwana to średnia
  \item $\textrm{Var}[X] = \mathbb{E}[(X - \mu)^2] = \mathbb{E}[X^2]
    - \mathbb{E}[X]^2 = \textrm{Cov}[X, X] = \sigma^2$
  \item $\sigma = \sqrt{\textrm{Var}[X]}$ - odchylenie standardowe
\end{itemize}

\subsection{Rozkład Normalny}

Rozkład normalny, lub rozkład Gaussa, jest rozkładem prawdopodobieństwa
zdefiniowanym na rzeczywistych liczbach rzeczywistych. Dla zmiennej losowej
$X$ o rozkładzie normalnym:
$$
X \sim \mathcal{N}(\mu, \sigma^2)
$$
gdzie $\mu$ to średnia, a $\sigma^2$ to wariancja.

Typowy estymator wariancji z próby:
$$
\hat{\sigma}^2(X) = \frac{1}{n} \sum_{i=1}^n (X_i - \overline{X})^2
\quad \mathbb{E}[\hat{\sigma}^2(X)] = \frac{n}{n-1} \textrm{Var}[X]
$$
jak widać, jest obciążony, albowiem $\mathbb{E}[\hat{\sigma}^2(X)]
\neq \textrm{Var}[X]$. Estymator nieobciążony wariancji to:
$$
S^2(X) = \frac{1}{n - 1} \sum_{i=1}^n (X_i - \overline{X})^2
\quad \mathbb{E}[S^2(X)] = \textrm{Var}[X]
$$

\subsection{Metoda najmniejszych kwadratów}

Jest to metoda estymowania parametrów modelu. Model ma postać $f(x, \theta)$,
gdzie $\theta = (\theta_1, \theta_2, \dots, \theta_m)$. Celem jest znalezienie
wartości $\theta$ dla której minimalizowane jest:
$$
r_i = y_i - f(x_i, \theta) \quad S = \sum_{i=1}^n r_i^2
$$

Rozwiązaniem problemu jest znalezienie takich $m$, dla których
gradient $S$ jest równy zero:
$$
\nabla_\theta S = 0
$$

\subsection{Metoda największej wiarygodności}

W tej metodzie mamy funkcję wiarygodności:
$$
L(\theta) = \prod_{i=1}^n f(y_i, \theta)
$$
gdzie $f(y_i, \theta)$ to prawdopodobieństwo wystąpienia wartości $y_i$ dla
parametrów $\theta$. Celem jest znalezienie takich $\theta$, dla których
$L(\theta)$ jest maksymalne.

\end{document}
